\documentclass[10pt,a4paper]{amsart}
\usepackage[margin=19mm]{geometry}
\usepackage{amsmath,amssymb,mathtools}
\usepackage[scr=rsfs,cal=euler]{mathalfa}
\usepackage{xfrac}
\usepackage[unicode,hidelinks,bookmarks=false]{hyperref}
\newcommand{\ie}{i.e.\ }
\newcommand{\cf}{cf.\ }
\newcommand{\resp}{respectively\ }
\newcommand{\Subsection}{\subsection}
\providecommand{\nobreakdash}{\nobreak\hbox{-}}
\setlength{\emergencystretch}{2em}
\multlinegap0pt
\usepackage{bm}
\usepackage{bbm}
\usepackage{dsfont}
\usepackage{xspace}
\usepackage{cancel}
\usepackage{mathtools}
\usepackage{amsthm}
\makeatletter
\@ifundefined{equation}{\newtheorem{equation}{Equation}}{}
\@ifundefined{definition}{\newtheorem{definition}[equation]{Definition}}{}
\@ifundefined{proposition}{\newtheorem{proposition}[equation]{Proposition}}{}
\@ifundefined{theorem}{\newtheorem{theorem}[equation]{Theorem}}{}
\makeatother
\newcommand{\C}{{\mathbb C}}
\newcommand{\abs}[1]{\left\vert#1\right\vert}
\newcommand{\cs}{{\mathcal S}}
\newcommand{\dbar}{\bar \partial}
\newcommand{\norm}[1]{\left\Vert#1\right\Vert}
\newcommand{\ol}{\overline}
\newcommand{\wt}{\widetilde}
\title[Projections onto $L^p$-Bergman spaces of Reinhardt Domains]{Projections onto $L^p$-Bergman spaces of Reinhardt Domains}
\author{Debraj Chakrabarti, Luke D. Edholm}
\date{Source: arXiv:2303.10005v3 (CC BY 4.0)}
\begin{document}
\maketitle

\noindent\emph{Source and licence.} Projections onto $L^p$-Bergman spaces of Reinhardt Domains. Version arXiv:2303.10005v3 (CC BY 4.0), distributed under the Creative Commons Attribution 4.0 International licence, as recorded in the arXiv licence metadata for this exact version and shown on the abstract page https://arxiv.org/abs/2303.10005v3. Full rights notice: licenses/arxiv-2303.10005-CC-BY-4.0.txt.\par\smallskip

\noindent\textit{Editorial scope.} Counted unit: the Proposition whose source label is \texttt{prop:MBP-is-a-basis-projection} in the article (source0.tex lines 788--797, source label \texttt{prop:MBP-is-a-basis-projection}), delivered together with the article's own complete original proof of it (source0.tex lines 798--843, one \texttt{proof} environment of 46 lines). No mathematics is written, repaired or re-derived: statement and proof are the article's own text line for line. Required context retained uncounted: T:monomials-form-a-Schauder-basis (lines 537--543); D:def-of-MBP-1 (lines 556--562); E:MBK-alt-form (lines 580--582); thm-MBK1 (lines 584--586); eq-kgalpha (lines 783--785). Every definition, display and label of those lines is retained unchanged. Omitted from this package and not redistributed: 1--536, 544--555, 563--579, 583--583, 587--782, 786--787, 844--1927. Nothing inside a retained excerpt is omitted or altered: every retained body is delivered exactly as the article's lines are written, so no omission receipt of any kind accompanies this package. Citations: the retained excerpts carry no citation command at all, so no bibliography entry of the article is transcribed and no reference is redistributed. Reference adaptation: none was needed and none was made; every reference inside the delivered unit resolves inside it, so no unresolved reference or citation is produced. Printed slips kept literal: no printed word, symbol, hypothesis or number of the counted statement or of its proof was altered. Layout and class: the article's own class is \texttt{amsart} with packages including inputenc, amsmath, amsthm, amssymb, amsfonts, mathrsfs, stmaryrd, amsxtra, epsfig, verbatim, enumerate, xcolor, enumitem, bigints; the wrapper is the lane's pinned \texttt{amsart} editorial wrapper at 10pt on A4, which restates the theorem-like environments the retained text needs and adds the article's own macro definitions that the retained text needs (\texttt{\textbackslash C}, \texttt{\textbackslash abs}, \texttt{\textbackslash cs}, \texttt{\textbackslash dbar}, \texttt{\textbackslash norm}, \texttt{\textbackslash ol}, \texttt{\textbackslash wt}), with extra wrapper packages (bm, bbm, dsfont, xspace, cancel, mathtools). The delivered numbering is the unit's own. Assets: no retained span contains a figure environment, a graphics-inclusion command, a PDF-page inclusion, a table or a TikZ picture; no image file is copied into this package, the package declares no asset dependency and no image file is redistributed with it. Licence: version arXiv:2303.10005v3 of this article is distributed under the Creative Commons Attribution 4.0 International licence, as recorded in the arXiv licence metadata for this exact version and shown on the abstract page https://arxiv.org/abs/2303.10005v3. A full rights notice is delivered at licenses/arxiv-2303.10005-CC-BY-4.0.txt. Characters: every retained excerpt is pure ASCII, so the delivered engine, XeLaTeX with its default Latin Modern text font, reports no missing character.
\begin{theorem}\label{T:monomials-form-a-Schauder-basis}
The collection of Laurent monomials $\{e_\alpha: \alpha \in \cs_p(\Omega,\lambda) \}$ forms a Banach-space basis of $A^p(\Omega,\lambda)$. 
The functionals dual to this basis are the coefficient functionals $\{a_\alpha: \alpha \in \cs_p(\Omega,\lambda) \}$, and the norm of $a_\alpha:A^p(\Omega, \lambda) \to \C$ is given by
\begin{equation}\label{eq-aalphanorm} 
\norm{{a_\alpha}}_{A^p(\Omega, \lambda)'} = \frac{1}{\norm{e_\alpha}_{p,\lambda}}.
\end{equation}
\end{theorem}

\begin{definition}\label{D:def-of-MBP-1}
A bounded linear projection $\bm{P}_{p,\lambda}^\Omega$ from $L^p(\Omega,\lambda)$ onto $A^p(\Omega,\lambda)$ is called the \emph{Monomial Basis Projection} of $A^p(\Omega,\lambda)$, if for $f \in L^p(\Omega,\lambda)$ it admits the series representation convergent in the norm of $L^p(\Omega,\lambda)$ given by
\begin{equation}\label{E:def-of-MBP-2}
\bm{P}^\Omega_{p,\lambda}(f) = \lim_{N\to \infty} \sum_{\substack{\abs{\alpha}_\infty\leq N \\ \alpha\in \cs_p(\Omega,\lambda) }}\wt{a}_\alpha(f) e_\alpha,
\end{equation}
where $\wt{a}_\alpha:L^p(\Omega,\lambda) \to \C$ is the unique Hahn-Banach extension of the coefficient functional $a_\alpha:A^p(\Omega,\lambda) \to \C$.
\end{definition}

\begin{equation}\label{E:MBK-alt-form}
    K_{p,\lambda}^\Omega(z,w) =\sum_{\alpha\in \mathcal{S}_p(\Omega, \lambda) }\frac{e_\alpha(z) \ol{e_\alpha(w)}\abs{e_\alpha(w)}^{p-2}} {\norm{e_\alpha}_{p,\lambda}^p}.
\end{equation}

\begin{theorem}\label{thm-MBK1}
Let $\Omega$ be a \emph{pseudoconvex}  Reinhardt domain in $\C^n$ and $\lambda$ be an admissible  multi-radial weight function on $\Omega$. The series \eqref{E:MBK-alt-form} defining $K^\Omega_{p,\lambda}(z,w)$ converges locally normally on  $\Omega \times \Omega$.  
\end{theorem}

\begin{equation}\label{eq-kgalpha}
K_{p,\lambda}^\Omega(z,w)=\sum_{\alpha\in \cs_p(\Omega, \lambda)}e_\alpha(z) \ol{g_\alpha(w)}. 
\end{equation}

\begin{proposition}\label{prop:MBP-is-a-basis-projection}
Define an integral operator on $C_c(\Omega)$ by
\begin{equation}\label{eq-qdef}
\bm{Q}f(z) = \int_\Omega  K_{p,\lambda}^\Omega(z,w) f(w)\lambda(w) dV(w), \qquad f\in C_c(\Omega).
\end{equation}
The MBP of $A^p(\Omega,\lambda)$ exists if and only if $\bm{Q}$ satisfies a weighted $L^p$-estimate, i.e., there is a constant $C>0$ such that for each $f\in C_c(\Omega)$ we have the inequality
\begin{equation}\label{eq-lpest}
\norm{\bm{Q} f}_{p,\lambda}\leq C \norm{f}_{p,\lambda}.
\end{equation}
\end{proposition}

\begin{proof} Recall that  $\Omega \subset \C^n$ is a pseudoconvex Reinhardt domain and $\lambda$ is an admissible multi-radial weight. 
The function $K_{p,\lambda}^\Omega$ is continuous on $\Omega\times \Omega$ by Theorem~\ref{thm-MBK1}, 
so the integral in \eqref{eq-qdef} exists for each $z\in \Omega$. 
Since the function $z \mapsto K_{p,\lambda}^\Omega(z, w)$ is holomorphic for each $w \in \Omega$, $\bm{Q}f$ is holomorphic for $f\in C_c(\Omega)$, for instance, by applying Morera's theorem in each variable, or equivalently, by applying $\dbar$ to both sides.

Let $f\in C_c(\Omega)$. 
Since the series for $K_{p,\lambda}^\Omega$ converges absolutely and uniformly on the compact subset $\{z\}\times \mathrm{supp}(f) \subset \Omega\times \Omega$, equation \eqref{eq-kgalpha} gives
\begin{align}
\bm{Q}f(z) &= \int_\Omega \bigg(\sum_{\alpha \in \cs_p(\Omega,\lambda)} e_\alpha(z)\ol{g_\alpha(w)}\bigg)f(w)\lambda(w)\,dV(w) \notag \\
&=\sum_{\alpha \in \cs_p(\Omega,\lambda)} \left( \int_\Omega f(w)  \ol{g_\alpha(w)}\lambda(w)\,dV(w)\right)e_\alpha(z)
=\sum_{\alpha \in \cs_p(\Omega,\lambda)} \wt{a}_\alpha(f) e_\alpha(z). \label{eq-qflaurent}
\end{align}
The series \eqref{eq-qflaurent} converges unconditionally and is the Laurent series of the holomorphic function $\bm{Q}f$. It is therefore uniformly convergent for $z$ in compact subsets of $\Omega$.

Suppose now that the MBP $\bm{P}_{p,\lambda}^\Omega:L^p(\Omega, \lambda)\to A^p(\Omega, \lambda) $ exists, which by Definition~\ref{D:def-of-MBP-1} is a bounded, surjective, linear projection given by the following limit of partial sums, convergent in $A^p(\Omega,\lambda)$:
\begin{equation}\label{eq-sqpartialsums} 
\bm{P}_{p,\lambda}^\Omega f = \lim_{N\to \infty} \sum_{\substack{\abs{\alpha}_\infty\leq N\\\alpha\in \cs_p(\Omega,\lambda) }}\wt{a}_\alpha(f) e_\alpha, \qquad f \in L^p(\Omega,\lambda).
\end{equation}
Since convergence in $A^p(\Omega,\lambda)$ implies uniform convergence on compact subsets, it follows that for $f \in C_c(\Omega)$, $\bm{Q}f = \bm{P}_{p,\lambda}^\Omega f$. 
Therefore $\bm{Q}$ satisfies $L^p$-estimates, i.e. \eqref{eq-lpest} holds.

Conversely, suppose that \eqref{eq-lpest} holds.  
Then $\bm{Q}$ can then be extended by continuity to an operator $\wt{\bm{Q}}$ on $L^p(\Omega, \lambda)$ with the same norm. 
We claim that $\wt{\bm{Q}}$ is the MBP.

If $f\in L^p(\Omega, \lambda)$, we can find a sequence $\{f_j\}\subset C_c(\Omega)$ such that $f_j\to f$ in $L^p(\Omega,\lambda)$. 
Each $\bm{Q}f_j \in A^p(\Omega, \lambda)$ and (by definition) $\bm{Q}f_j\to \wt{\bm{Q}}f$ in  $L^p(\Omega, \lambda)$. 
But this implies $\bm{Q}f_j\to \wt{\bm{Q}}f$ uniformly on compact subsets, so the limit  $\wt{\bm{Q}}f$ is holomorphic, and thus the range of $\wt{\bm{Q}}$ is contained in $A^p(\Omega, \lambda)$. 
A direct computation now shows $\wt{\bm{Q}}e_\alpha = e_\alpha$ for $\alpha \in \cs_p(\Omega,\lambda)$, and it follows that $\wt{\bm{Q}}$ is a surjective projection from $L^p(\Omega,\lambda)$ to $A^p(\Omega, \lambda)$. 

If $f\in C_c(\Omega) $, then ${\bm{Q}}f = \wt{\bm{Q}}f \in A^p(\Omega, \lambda)$ and by Theorem~\ref{T:monomials-form-a-Schauder-basis} the Laurent series expansion of $\wt{\bm{Q}}f$ given by \eqref{eq-qflaurent} converges (as a sequence of square partial sums) in $A^p(\Omega,\lambda)$:
\begin{equation}\label{eq-qrearrange}
\wt{\bm{Q}}f = \lim_{N\to \infty} \sum_{\substack{\abs{\alpha}_\infty\leq N\\\alpha\in \cs_p(\Omega,\lambda) }}\wt{a}_\alpha(f) e_\alpha. 
\end{equation}
For a general $g \in L^p(\Omega, \lambda)$, $\wt{\bm{Q}}g \in A^p(\Omega,\lambda)$ and so again by Theorem~\ref{T:monomials-form-a-Schauder-basis},
\begin{equation}\label{eq-wtbmqg} 
\wt{\bm{Q}}g = \lim_{N\to \infty} \sum_{\substack{\abs{\alpha}_\infty\leq N\\\alpha\in \cs_p(\Omega,\lambda)}} {a}_\alpha(\wt{\bm{Q}}g) e_\alpha.
\end{equation} 
It follows that on $C_c(\Omega)$ we have the identity $a_\alpha\circ{\bm{Q}}= \wt{a}_\alpha$. 
This relation extends by continuity to give  $a_\alpha\circ{\wt{\bm{Q}}}= \wt{a}_\alpha$ as functionals on $L^p(\Omega, \lambda)$. 
Then \eqref{eq-wtbmqg} becomes
\[
\wt{\bm{Q}}g = \lim_{N\to \infty} \sum_{\substack{\abs{\alpha}_\infty\leq N\\\alpha\in \cs_p(\Omega,\lambda) }}\wt{a}_\alpha(g) e_\alpha.
\]
In other words, $\wt{\bm{Q}}$ is the MBP, as we wanted to show. 
\end{proof}

\end{document}
